% Tutorial set: 04
% Source: T4.pdf, Tutorial 4 (Image Edge Processing), Questions 1--3
% Reference: T4_Reference solutions.pptx, used only for answer checking
% Session date: 2026-09-15
% Human verified: Yes

\clearpage
\exercisegroup{SC4061-TUT04}{Tutorial 04：Image Edge Processing}
\label{tutorial:SC4061-TUT04}

\exercisemeta
  {SC4061-TUT04}
  {2026-09-15}
  {T4.pdf，Questions 1--3；T4\_Reference solutions.pptx 仅用于核对}
  {三题均已作答、讨论并完成订正}
  {已确认（2026-09-15）}

\exercisequestion{SC4061-TUT04-Q01}{Sobel Mask 的 Separable Implementation}

\textbf{对应讲义：}
\relatedtopic{SC4061-L05-T05}{Prewitt 与 Sobel Gradient Filters}

\subsubsection*{题目}

证明 Sobel masks（Sobel 卷积核）可以由一次一维 differencing（差分）和一次正交方向的一维 smoothing（平滑）依次实现：差分核为 $[-1,0,1]^{\mathsf T}$ 或其水平版本，平滑核为 $[1,2,1]$ 或其垂直版本。

\begin{checkedsolution}
令水平差分与垂直平滑分别为
\[
  d_x=\begin{bmatrix}-1&0&1\end{bmatrix},\qquad
  s_y=\begin{bmatrix}1&2&1\end{bmatrix}^{\mathsf T}.
\]
二者的 outer product（外积）为
\[
  K_x=s_y d_x
  =\begin{bmatrix}
    -1&0&1\\
    -2&0&2\\
    -1&0&1
  \end{bmatrix}.
\]
$K_x$ 计算沿 $x$ 的 intensity difference（强度差），因此响应 vertical edges（竖直边缘）。交换两个方向可得
\[
  K_y=d_y s_x
  =\begin{bmatrix}
    -1&-2&-1\\
    0&0&0\\
    1&2&1
  \end{bmatrix},
\]
它计算沿 $y$ 的差分并响应 horizontal edges（水平边缘）。

由 convolution（卷积）的结合律，二维 Sobel filtering 可写成
\[
  f*K_x=(f*d_x)*s_y=(f*s_y)*d_x.
\]
因此先差分再平滑或先平滑再差分都得到同一二维结果。若软件实际执行 cross-correlation（互相关）而没有翻转 kernel，输出可能整体变号；这会改变 gradient orientation（梯度朝向）$180^\circ$，但不会改变 gradient magnitude 或 edge location。
\end{checkedsolution}

\exercisequestion{SC4061-TUT04-Q02}{LoG Kernel、Filtering Response 与 Zero Crossing}

\textbf{对应讲义：}
\relatedtopic{SC4061-L05-T06}{Laplacian of Gaussian 与 Zero Crossing}

\subsubsection*{题目}

令
\[
  h(r)=-\exp\!\left(-\frac{r^2}{2\sigma^2}\right),
  \qquad r^2=x^2+y^2.
\]
（a）详细推导 $\nabla^2h(r)$。（b）使用离散 LoG kernel
\[
  L=\begin{bmatrix}
    0&0&-1&0&0\\
    0&-1&-2&-1&0\\
    -1&-2&16&-2&-1\\
    0&-1&-2&-1&0\\
    0&0&-1&0&0
  \end{bmatrix}
\]
对两幅 $12\times13$ 图像做 valid filtering（不补零，只保留 kernel 完全落入图像的位置），并比较结果。两幅图像的 12 行分别重复以下 row profiles：
\[
  x_{\mathrm{step}}=
  [0,0,0,0,0,0,1,1,1,1,1,1,1],
\]
\[
  x_{\mathrm{ramp}}=
  [0,0,0,0,0,0,0.5,1,1,1,1,1,1].
\]

\begin{checkedsolution}
记 $q=(x^2+y^2)/(2\sigma^2)$。链式法则给出
\[
  \frac{\partial h}{\partial x}
  =\frac{x}{\sigma^2}e^{-q},\qquad
  \frac{\partial^2h}{\partial x^2}
  =-\frac{x^2-\sigma^2}{\sigma^4}e^{-q}.
\]
同理，
\[
  \frac{\partial^2h}{\partial y^2}
  =-\frac{y^2-\sigma^2}{\sigma^4}e^{-q}.
\]
所以
\[
  \boxed{
  \nabla^2h(r)
  =-\frac{r^2-2\sigma^2}{\sigma^4}
  \exp\!\left(-\frac{r^2}{2\sigma^2}\right)}.
\]
题目采用的是 negative unnormalised Gaussian（负的未归一化高斯函数），所以 kernel 中心为正。若从通常的 positive normalised Gaussian
\[
  G_\sigma(r)=\frac{1}{2\pi\sigma^2}
  e^{-r^2/(2\sigma^2)}
\]
出发，则
\[
  \nabla^2G_\sigma(r)
  =\frac{r^2-2\sigma^2}{2\pi\sigma^6}
  e^{-r^2/(2\sigma^2)},
\]
与题目 convention（约定）整体异号并多一个归一化系数；二者的 zero crossing 都位于 $r=\sqrt2\sigma$。

因为图像的每一行都相同，二维 filtering 可约化为一维。将 $L$ 的每列求和，得到 effective kernel（等效卷积核）
\[
  \ell=[-1,-4,10,-4,-1].
\]
输入高度与宽度分别从 $12,13$ 变成
\[
  12-5+1=8,\qquad 13-5+1=9.
\]
每个输出均有 8 个相同行，其 row profiles 为
\[
  \boxed{y_{\mathrm{step}}
  =[0,0,-1,-5,5,1,0,0,0]},
\]
\[
  \boxed{y_{\mathrm{ramp}}
  =[0,0,-0.5,-3,0,3,0.5,0,0]}.
\]
Reference solutions 的渐变响应末尾漏写了一个 $0$；由 $13-5+1=9$ 可确认完整输出必须有 9 列。

Abrupt step（突变阶跃）产生较强的正负 double-lobe response（双瓣响应），zero crossing 位于 $-5$ 与 $5$ 之间。加入 $0.5$ 的过渡后，响应峰值降至 $\pm3$，并在边缘中心出现一个精确的零。LoG 以 sign change（符号变化）定位边缘；整体改变 kernel sign 不改变 zero-crossing location。
\end{checkedsolution}

\exercisequestion{SC4061-TUT04-Q03}{Hough Line Geometry 与 Parameter-space Voting}

\textbf{对应讲义：}
\relatedtopic{SC4061-L06-T01}{Edgel Linking 与 Hough Line Geometry}；
\relatedtopic{SC4061-L06-T02}{Image--Hough Duality、Voting 与 Generalization}

\subsubsection*{题目}

直线可写为
\[
  \rho=x\cos\theta+y\sin\theta.
\]
（a）从几何上说明 $\rho$ 与 $\theta$；（b）在可行时将直线写成 $y=mx+c$；（c）固定图像点 $(x,y)$，证明它在 $(\theta,\rho)$ parameter space（参数空间）中对应 sinusoid（正弦曲线）$\rho=K\cos\alpha$；（d）解释两条以及多条 sinusoids 的交点分别对应什么图像结构。

\begin{checkedsolution}
直线的单位法向量为
\[
  \mathbf n=(\cos\theta,\sin\theta)^{\mathsf T}.
\]
$\theta$ 是 $\mathbf n$ 与 $x$ 轴的夹角，$\rho$ 是原点到直线的 signed distance（有符号距离），而距离原点最近的直线点为 $\rho\mathbf n$。

当 $\sin\theta\ne0$ 时，
\[
  y=-\frac{\cos\theta}{\sin\theta}x
  +\frac{\rho}{\sin\theta}
  =-\cot\theta\,x+\rho\csc\theta.
\]
因此 $m=-\cot\theta$，$c=\rho\csc\theta$。当 $\sin\theta=0$ 时，直线为 $x=\rho/\cos\theta$，即 vertical line（竖直线），不能写成具有有限 slope 的 $y=mx+c$。Normal form（法线式）正是通过 $\rho,\theta$ 避免 slope singularity（斜率奇点）。

对固定点 $(x,y)$，令
\[
  R=\sqrt{x^2+y^2},\qquad
  \phi=\operatorname{atan2}(y,x),\qquad
  x=R\cos\phi,\quad y=R\sin\phi.
\]
代入直线方程：
\begin{align*}
  \rho(\theta)
  &=R\cos\phi\cos\theta+R\sin\phi\sin\theta\\
  &=\boxed{R\cos(\theta-\phi)}.
\end{align*}
所以题目中的 $K=R$，而 $\alpha$ 可取 $\theta-\phi$；由于 cosine（余弦）是偶函数，写成 $R\cos(\phi-\theta)$ 也完全等价。Reference solutions 使用 $\tan^{-1}(y/x)$ 表示 $\phi$，但严格实现时应使用 $\operatorname{atan2}(y,x)$，否则不能正确区分象限，且 $x=0$ 时会发生除零。

每个 edgel（边缘像素）在 Hough space 中给所有经过它的直线投票，因此形成一条 sinusoid。两条 sinusoids 的交点表示两个图像点共同确定的一条直线；多条 sinusoids 在同一 $(\rho,\theta)$ 附近相交，表示多个 edgels 共线，并在 accumulator（累加器）中形成较强 peak。离散实现中受 parameter quantisation（参数量化）和噪声影响，投票通常形成邻近 bins 的峰簇，而不是理想的精确交点。

\begin{figure}[htbp]
  \centering
  \resizebox{0.94\textwidth}{!}{\begin{tikzpicture}[font=\small,>=Latex]
  \begin{scope}[xshift=0cm]
    \draw[->,thick] (-0.3,0) -- (3.4,0) node[right] {$x$};
    \draw[->,thick] (0,-1.3) -- (0,1.7) node[above] {$y$};
    \draw[very thick,noteBlue] (1,-1.15) -- (1,1.45);
    \foreach \yy/\cc in {-0.7/ntuRed,0/orange!85!black,0.7/teal!75!black} {
      \fill[\cc] (1,\yy) circle (3.2pt);
    }
    \node[align=center] at (1.65,-1.55) {image space\\共线 edgels};
  \end{scope}

  \draw[->,very thick,gray!65] (3.7,0.2) -- (5.1,0.2)
    node[midway,above,align=center] {vote};

  \begin{scope}[xshift=5.7cm,x=0.035cm,y=1.05cm]
    \draw[->,thick] (-85,0) -- (88,0) node[right] {$\theta$};
    \draw[->,thick] (0,-1.35) -- (0,1.7) node[above] {$\rho$};
    \draw[ntuRed,very thick,samples=90,domain=-80:80,smooth]
      plot (\x,{cos(\x)-0.7*sin(\x)});
    \draw[orange!85!black,very thick,samples=90,domain=-80:80,smooth]
      plot (\x,{cos(\x)});
    \draw[teal!75!black,very thick,samples=90,domain=-80:80,smooth]
      plot (\x,{cos(\x)+0.7*sin(\x)});
    \fill[noteBlue] (0,1) circle (3.5pt);
    \draw[dashed,noteBlue] (0,0) -- (0,1);
    \node[anchor=west,noteBlue] at (5,1.18) {peak $(\theta^*,\rho^*)$};
    \node[align=center] at (0,-1.65) {Hough space\\sinusoids intersect};
  \end{scope}
\end{tikzpicture}
}
  \caption{Q3 的 image--Hough duality：image space 中的共线 edgels 在 Hough space 中产生共同参数峰。图为依据课程内容绘制的原创示意图。}
  \label{fig:sc4061-tut04-hough-voting}
\end{figure}
\end{checkedsolution}

\subsubsection*{本次 Tutorial 收口}

Q1 利用 separability（可分离性）把二维 Sobel convolution 拆为差分与平滑；Q2 连接 continuous LoG derivation、discrete filtering response 与 zero crossing，并用输出尺寸发现参考答案漏项；Q3 在 image space 与 Hough space 之间切换变量角色，以 signed normal form 统一表示普通直线和竖直线。复习时应优先保留 kernel direction、sign convention、valid-output size、$\sin\theta=0$ 特例与 $\operatorname{atan2}$ 的象限信息。
